% alt: A staircase of boxes above a wide box Omega with five physical legs: P_4 sits on legs 1-4, P_3 on legs 1-3 above it, P_2 on legs 1-2, and P_1 on leg 1 at the top; the five legs leave the top of the picture labelled 1 to 5, and the whole picture equals Phi.
% Section 7, Proposition prop:VC (part2a_mps_information.tex, lines 25-47): Phi = (Pi_1 (x) I)(Pi_2 (x) I) ... (Pi_{L-1} (x) I) Omega with P_k = Pi_k (x) I as in the proof, drawn for L = 5; Pi_k is the projection of H_{[1,k]} onto the span of the D leading eigenvectors of rho_{[1,k]}, so each Pi_k has rank D.
% Ink-to-index: the wide bottom box is the unit vector Omega, its five legs the sites 1..5 of the chain from left to right; the box P_k is Pi_k (x) I and covers the legs 1..k (the identity factor on [k+1,L] is not drawn, those legs pass beside the box); P_4 is nearest Omega because P_{L-1} acts first. Every leg of Omega is contracted with the lowest box covering its site, every output of P_k with P_{k-1}; the five legs at the top are the open physical indices, labelled by their site. The vector P_k ... P_{L-1} Omega lies in ran Pi_k (x) H_{[k+1,L]}, so its Schmidt rank across the cut k is at most D, and the later factors P_1, ..., P_{k-1} act on [1,k-1] and do not increase it; hence Phi is a matrix product state of bond dimension at most D.
\[
  \Phi\;=\;
  \begin{tenkz}[rows={wire,wire,wire,wire,wire}, cols=5, bonds=none]
    \tn[at=(5,1), skin=box, wide=5, name=o,
        ports={90@1:physical, 90@2:physical, 90@3:physical, 90@4:physical, 90@5:physical}]{\Omega}
    \tn[at=(4,1), skin=box, wide=4, name=p4,
        ports={270@1:physical, 270@2:physical, 270@3:physical, 270@4:physical,
               90@1:physical, 90@2:physical, 90@3:physical, 90@4:physical}]{P_4}
    \tn[at=(3,1), skin=box, wide=3, name=p3,
        ports={270@1:physical, 270@2:physical, 270@3:physical,
               90@1:physical, 90@2:physical, 90@3:physical}]{P_3}
    \tn[at=(2,1), skin=box, wide=2, name=p2,
        ports={270@1:physical, 270@2:physical, 90@1:physical, 90@2:physical}]{P_2}
    \tn[at=(1,1), skin=box, name=p1,
        ports={270:physical, 90:physical:$1$}]{P_1}
    \tn[at=(1,2), skin=none, name=t2, ports={270:physical, 90:physical:$2$}]{}
    \tn[at=(1,3), skin=none, name=t3, ports={270:physical, 90:physical:$3$}]{}
    \tn[at=(1,4), skin=none, name=t4, ports={270:physical, 90:physical:$4$}]{}
    \tn[at=(1,5), skin=none, name=t5, ports={270:physical, 90:physical:$5$}]{}
    \tnwire{o.90@1}{p4.270@1} \tnwire{o.90@2}{p4.270@2}
    \tnwire{o.90@3}{p4.270@3} \tnwire{o.90@4}{p4.270@4}
    \tnwire{p4.90@1}{p3.270@1} \tnwire{p4.90@2}{p3.270@2} \tnwire{p4.90@3}{p3.270@3}
    \tnwire{p3.90@1}{p2.270@1} \tnwire{p3.90@2}{p2.270@2}
    \tnwire{p2.90@1}{p1.270}
    \tnwire{p2.90@2}{t2.270}
    \tnwire{p3.90@3}{t3.270}
    \tnwire{p4.90@4}{t4.270}
    \tnwire{o.90@5}{t5.270}
  \end{tenkz}
\]
